% !TeX program = pdflatex
\documentclass[a4paper,11pt]{article}
\usepackage[T1]{fontenc}
\usepackage[utf8]{inputenc}
\usepackage{lmodern}
\usepackage[english]{babel}
\usepackage[margin=1in]{geometry}
\usepackage{graphicx}
\usepackage{caption}
\usepackage{hyperref}
\hypersetup{colorlinks=true,linkcolor=black,citecolor=black,urlcolor=black,
  pdftitle={Milan Housing Project Report},pdfauthor={Mattia Faini}}
% The build script extracts this artwork from the unchanged Italian PDF.
\graphicspath{{.build/report/figures/}}
\title{Milan Housing Project Report}
\author{Mattia Faini}
\date{}

\begin{document}
\maketitle
\begin{center}
\small English translation prepared on 27 September 2026. Original figure artwork
is retained; English keys accompany the Italian labels.
\end{center}

\section{Data cleaning, missing-value imputation and predictor transformations}
The \texttt{training} and \texttt{test} sets contained some missing or inconsistent
values. Transformations were also considered appropriate for several predictors.
Each dataset field was treated as follows:

\begin{itemize}
  \item Missing values of \texttt{bathrooms\_number} were imputed using cumulative
  logistic regression, with \texttt{square\_meters} as the predictor.

  \item \texttt{rooms\_number} had no missing values. The value \texttt{5+} was
  converted to \texttt{6}.

  \item Some \texttt{square\_meters} values were incompatible with
  \texttt{rooms\_number}. They were replaced by predictions
  from a linear model with \texttt{rooms\_number} as the predictor. The variable
  \texttt{sqm\_1}, equal to the reciprocal of the area in square metres, was also
  created.

  \item A new response variable, \texttt{sqm\_price}, was created from
  \texttt{selling\_price}. It is the logarithm of the price per square metre and
  is used to fit the final model and calculate predictions.

  \item Missing values of \texttt{year\_of\_construction} were imputed using the
  mean of the other values. An indicator, \texttt{new}, was also created: it equals
  \texttt{1} if the property was built after 2010.

  \item Missing values of \texttt{total\_floors\_in\_building} were imputed using
  a linear model with \texttt{year\_of\_construction} as the predictor.

  \item \texttt{car\_parking} was split into \texttt{shared\_parking} and
  \texttt{garage\_box}, which represent parking spaces in garages and communal
  areas, respectively.

  \item \texttt{availability} was converted to an indicator equal to \texttt{0}
  if the property is not immediately available and \texttt{1} if it is.

  \item \texttt{floor} was converted to a numeric variable. The values
  \texttt{ground}, \texttt{mezzanine} and \texttt{basement} were set to zero, and
  corresponding indicators were created. The variable \texttt{ratio}, the ratio
  of \texttt{floor} to \texttt{total\_floors\_in\_building}, was also created.

  \item \texttt{energy\_effciency\_class} was converted to numeric values.
  Missing values were predicted using cumulative logistic regression.

  \item An indicator was created for each feature in \texttt{other\_features},
  including an \texttt{NA} indicator for properties with no reported features.

  \item \texttt{lift} was converted to an indicator equal to \texttt{0} if the
  property has a lift and \texttt{1} if it does not. Missing values were predicted
  using logistic regression. The variable \texttt{floor\_lift}, the product of
  \texttt{floor} and \texttt{lift}, was also created to identify properties on
  higher floors without a lift.

  \item \texttt{condominium\_fees} was converted to numeric values. It was divided
  by \texttt{square\_meters}, and missing values were imputed using the mean.

  \item Missing values of \texttt{heating\_centralized} and \texttt{conditions}
  were imputed using the mode.

  \item The values of \texttt{zone} were geolocated using the Google Maps API
  and imported into the geographical software \textsf{QGIS}. Geolocation errors
  were corrected. As shown in \autoref{fig:map}, a legend was added to show the
  number of observations for each zone in the \texttt{training} set. Zones with
  fewer than 3 observations in \texttt{training} were merged with neighbouring
  zones.
\end{itemize}

\begin{figure}[!htbp]
  \centering
  \includegraphics[width=\textwidth]{artwork-000.png}
  \caption{Map of Milan showing the zones present in the datasets. The legend
  gives the number of \texttt{training} observations in each zone.}
  \label{fig:map}
  \par\smallskip
  {\small\raggedright\textbf{English key.} ``n\textsuperscript{o} osservazioni in
  training.csv'' means ``number of observations in training.csv''. The categories
  are 0--1, 2--3, 4--10, 11--20 and 21--300. Neighbourhood and place names are
  retained as proper names.\par}
\end{figure}

\clearpage
\section{The model}
Because the loss to be minimised was absolute error (MAE), a quantile regression
model was chosen. In this project, implementing dimension-reduction or shrinkage
techniques such as lasso, group lasso, ridge and PCA with quantile regression
proved difficult. Trials of these techniques on the dataset did not yield satisfactory results.
Quantile regression was therefore preferred despite its greater computational
complexity.

\subsection{Attempt to reduce the dimensionality of the zones}
An attempt was made to merge zones using k-means on their coordinates. The results
were assessed on a validation test subset drawn randomly from \texttt{training},
while dividing the observations from each zone proportionally between the training
and validation test subsets. A series of regressions was fitted using all predictors, with the
zones merged into progressively fewer groups, in a procedure similar to backward
regression.

As \autoref{fig:kmeans} shows, the error increases almost monotonically as the
number of zones decreases. All zones were therefore retained in the subsequent
steps.

\begin{figure}[!htbp]
  \centering
  \includegraphics[width=\textwidth]{artwork-002.jpg}
  \caption{Mean absolute error on the test set as a function of the number of
  distinct zones included. The number of zones was reduced by merging them using
  k-means on their geographical coordinates. The vertical line marks the point
  with the lowest MAE.}
  \label{fig:kmeans}
  \par\smallskip
  {\small\raggedright\textbf{English key.} ``MAE in funzione del numero di zone
  incluse'' means ``MAE as a function of the number of zones included''.
  ``Numero di zone incluse'' means ``Number of zones included''.
  MAE means mean absolute error.\par}
\end{figure}

\subsection{Final predictor selection}
Backward regression was tried as a way to reduce dimensionality, removing the
predictor with the highest p-value at each step. However, the results varied
substantially with the choice of training and test sets. Fifty different models
were therefore selected using backward regression, and each was evaluated on 50
different training and test splits (\autoref{fig:hist}). The model selected for
the final predictions had the lowest MAE and an acceptable variance.

A limitation of this procedure is that the data used for the final evaluation had
already been used to select the 50 candidate models. The model's MAE estimate is
also biased because that same estimate was used to choose the model.

\begin{figure}[!htbp]
  \centering
  \includegraphics[width=\textwidth]{artwork-003.jpg}
  \caption{Distributions of the number of predictors, MAE and standard deviation
  of MAE for 50 models obtained by backward regression on 50 training and test
  splits and evaluated on another 50 training and test splits. The vertical lines
  indicate the values for the selected final model. The zone indicators are
  counted as one predictor, as are the indicators for \texttt{conditions}.}
  \label{fig:hist}
  \par\smallskip
  {\small\raggedright\textbf{English key.} ``MAE medio'' means ``Mean MAE''.
  ``MAE Std'' and ``Deviazione standard del MAE'' mean ``Standard deviation of
  MAE''. ``Numero di predittori'' means ``Number of predictors''.
  ``Frequency'' is already in English.\par}
\end{figure}
\end{document}
